\begin{abstract}
This paper reports the design, firmware implementation, object-oriented
architecture audit, and dummy-cell performance evaluation of HunStat2, our
own low-cost, open-hardware potentiostat built around the Analog Devices
AD5940 electrochemical analog front end (AFE) and a Seeed Studio XIAO
RP2040 microcontroller, built directly from the AD5940 datasheet, standard
three-electrode potentiostat design practice, and the AD5940
evaluation-kit reference firmware. The firmware is organized as a partial
migration from a monolithic, procedural Arduino sketch toward a set of C++
classes separating command processing, parameter storage, hardware
configuration, and individual electrochemical methods (chronoamperometry,
CA; cyclic voltammetry, CV; square-wave voltammetry, SWV; differential
pulse voltammetry, DPV; electrochemical impedance spectroscopy, EIS; and
open-circuit potential, OCP). We report an honest, evidence-based audit of
this structure's object-oriented maturity: of eight classes, only two
exhibit non-trivial encapsulation, no inheritance or polymorphism is used
anywhere in the codebase, and three method classes contained
verbatim-duplicated measurement-configuration code -- duplication we show
was not a cosmetic concern, since a real correctness bug in that shared
code had to be found and fixed three times over, once per copy. Direct,
file-by-file comparison of the project's CA/SWV/DPV firmware against
Analog Devices' own AD5940 reference examples identified exactly that bug,
and a second, independent one, together explaining a previously-reported
CA/SWV data-quality failure: (1) the affected classes configured the
AD5940's High-Speed excitation/sensing loop, intended for AC/impedance
excitation, where every vendor DC-method reference example instead
configures the Low-Power loop; and (2) the analog-to-digital conversion
routine used the wrong zero-point convention for the chip's offset-binary
signed codes, turning small, physically sensible signals into large,
discontinuous, non-physical values. Both defects were corrected
identically across all three classes and are carried forward into the same
firmware baseline independently validated on a traceable dummy cell by two
companion manuscripts: that data confirms CV, CA, SWV,
and DPV all produce correctly-signed, signal-dependent, non-degenerate
output, and the companion metrology manuscript's cyclic-voltammetry
accuracy check ($R^2\geq0.999994$, $+0.428\,\%$ recovered-resistance error
against a traceable Randles dummy cell) is independent confirmation that
the corrected signal chain is quantitatively correct, not merely
self-consistent. This paper positions its contribution around a diagnosed,
evidence-based hardware-and-software architecture assessment together with
a working, hardware-validated measurement chain, rather than an
unsubstantiated claim of a fully mature object-oriented design; the
structural refactor that the audit motivates is reported and separately
validated in a companion manuscript built specifically around it.
\end{abstract}

\begin{keyword}
Potentiostat \sep AD5940 \sep object-oriented design \sep embedded systems
\sep electrochemical instrumentation \sep firmware debugging \sep
reference-code comparison \sep dummy cell \sep open-source hardware
\end{keyword}
